%======================================================================
% CAPÍTULO 2: MARCO TEÓRICO - MOCKUP ESTILO REVISTA CIENTÍFICA
%======================================================================
\chapter{Marco Te\'orico}

%--------------------------------------------------------------------
% BANNER DE CAPÍTULO
%--------------------------------------------------------------------
\bannerCapitulo
    {Marco Te\'orico}
    {Fundamentos tecnol\'ogicos y arquitectura de sistemas empresariales}
    {02}

%--------------------------------------------------------------------
% ESTADO DEL ARTE
%--------------------------------------------------------------------
\section{Estado del Arte}

\dropcap{E}{l} middleware empresarial ha evolucionado significativamente desde sus or\'igenes hasta convertirse en plataformas complejas de integraci\'on. El mercado ofrece soluciones como IBM Integration Bus, MuleSoft, Apache Camel y Spring Integration.

Sin embargo, la mayor\'ia fueron dise\~nadas para mercados desarrollados, ignorando las necesidades espec\'ificas de regiones con infraestructura limitada y contextos regulatorios distintos. Esta brecha crea una oportunidad para soluciones localmente desarrolladas que aborden estos desaf\'ios particulares.

\begin{cajaISO}{ISO 27001:2022}
    La familia de normas ISO/IEC 27000 proporciona un marco sistem\'atico para gestionar la seguridad de la informaci\'on. Sandra Ecosystem implementa estos principios en todas las capas de su arquitectura, garantizando la confidencialidad, integridad y disponibilidad de los datos.
\end{cajaISO}

\begin{definicionMagazine}
\textbf{Estado del Arte}: Conjunto de tecnolog\'ias, m\'etodos y pr\'acticas m\'as avanzadas para resolver un problema espec\'ifico en un momento determinado.
\end{definicionMagazine}

%--------------------------------------------------------------------
% FUNDAMENTOS DE MIDDLEWARE  
%--------------------------------------------------------------------
\section{Fundamentos de Middleware}

El middleware es un software que act\'ua como puente entre aplicaciones y el sistema operativo, facilitando la comunicaci\'on e integraci\'on. Un Enterprise Service Bus (ESB) representa la evoluci\'on hacia una arquitectura basada en servicios.

\begin{cajadestacada}{Implementaci\'on en Sandra Server}
    Sandra Server implementa estos conceptos utilizando Go, proporcionando un ESB de alto rendimiento capaz de manejar miles de conexiones simult\'aneas con consumo m\'inimo de recursos.
\end{cajadestacada}

\begin{figure}[H]
    \centering
    \begin{tikzpicture}[
        node distance=2cm,
        app/.style={rectangle, draw, fill=tealpastel, rounded corners, minimum width=2cm, minimum height=1cm},
        mid/.style={rectangle, draw, fill=sandraMagAccent!20, rounded corners, minimum width=3cm, minimum height=1.2cm, font=\bfseries},
        os/.style={rectangle, draw, fill=secondary!20, rounded corners, minimum width=2cm, minimum height=1cm}
    ]
        \node[app] (app1) {App 1};
        \node[app] (app2) {App 2};
        \node[app, right=1cm of app2] (app3) {App N};
        
        \node[mid, below=1.5cm of app1] (esb) {Enterprise Service Bus};
        
        \node[os, below=1.5cm of esb] (so) {Sistema Operativo};
        
        \draw[->, thick] (app1) -- (esb);
        \draw[->, thick] (app2) -- (esb);
        \draw[->, thick] (app3) -- (esb);
        \draw[->, thick] (esb) -- (so);
    \end{tikzpicture}
    \caption{Arquitectura b\'asica de Middleware}
    \label{fig:middleware-arch}
\end{figure}

%--------------------------------------------------------------------
% ARQUITECTURA DE MICROSERVICIOS
%--------------------------------------------------------------------
\section{Arquitectura de Microservicios}

La arquitectura de microservicios construye aplicaciones como servicios peque\~nos, aut\'onomos y dbilmente acoplados. Cada microservicio es responsable de una funcionalidad espec\'ifica y puede desplegarse independientemente.

\begin{cajaInfo}
El ecosistema Sandra implementa esta arquitectura distribuida:
\end{cajaInfo}

\begin{listaTecnologias}
    \item \textbf{Server}: Middleware central de comunicaciones
    \item \textbf{Consola}: Interfaz administrativa web
    \item \textbf{SDC}: Aplicaci\'on de escritorio
    \item \textbf{Sentinel}: Procesamiento especializado de eventos
\end{listaTecnologias}

\begin{cajaAdvertencia}
    La migraci\'on de arquitecturas monol\'iticas a microservicios requiere una planificaci\'on cuidadosa. Se recomienda comenzar con servicios independientes y gradually desacoplar funcionalidad.
\end{cajaAdvertencia}

\begin{wrapfigure}{r}{0.4\textwidth}
    \centering
    \begin{tikzpicture}[
        node distance=0.8cm,
        ms/.style={circle, draw, fill=sandraMagAccent!30, minimum width=1.2cm, font=\bfseries\tiny}
    ]
        \node[ms] (s1) {Server};
        \node[ms] (s2) [right of=s1] {Consola};
        \node[ms] (s3) [below of=s1] {SDC};
        \node[ms] (s4) [below of=s2] {Sentinel};
        
        \draw[<->, thick] (s1) -- (s2);
        \draw[<->, thick] (s1) -- (s3);
        \draw[<->, thick] (s1) -- (s4);
    \end{tikzpicture}
    \caption{Ecosistema Sandra}
    \label{fig:sandra-ecosystem}
\end{wrapfigure}

%--------------------------------------------------------------------
% SEGURIDAD DE CONFIANZA CERO
%--------------------------------------------------------------------
\section{Seguridad de Confianza Cero}

\dropcap{E}{l} modelo Zero Trust elimina la confianza impl\'icita, verificando cada solicitud independientemente de su origen. Los principios fundamentales incluyen:

\begin{cajaISO}{ISO 27002:2022}
    Los controles de seguridad de la informaci\'on seg\'un ISO 27002 incluyen 93 medidas organizadas en cuatro categor\'ias: organizacionales, personales, f\'isicos y tecnol\'ogicos.
\end{cajaISO}

\begin{enumerate}
    \item \textbf{Verificaci\'on expl\'icita}: Autenticaci\'on multifactor para todos los accesos
    \item \textbf{M\'inimo privilegio}: Permisos m\'inimos necesarios para cada funci\'on
    \item \textbf{Verificaci\'on continua}: Monitoreo constante de sesiones y actividades
\end{enumerate}

Sandra Ecosystem implementa estos principios en todas las capas: cada solicitud de API es autenticada y autorizada, las comunicaciones usan TLS m\'utuo, y las sesiones tienen duraci\'on limitada.

%--------------------------------------------------------------------
% TECNOLOGÍAS EMERGENTES
%--------------------------------------------------------------------
\section{Tecnolog\'ias Emergentes}

El ecosistema utiliza tecnolog\'ias de \'ultima generaci\'on:

\begin{tabular}{m{3cm} m{10cm}}
    \textcolor{sandraMagAccent}{\faIcon{code}} \textbf{Go 1.24.3} & Alto rendimiento y concurrencia nativa para el middleware central \\[0.8em]
    \textcolor{sandraMagGold}{\faIcon{shield-alt}} \textbf{Rust 2021} & Seguridad de memoria y rendimiento para componentes cr\'iticos \\[0.8em]
    \textcolor{isoRed}{\faIcon{layer-group}} \textbf{Angular 16+} & Framework frontend moderno con soporte enterprise \\[0.8em]
    \textcolor{sandraMagDark}{\faIcon{desktop}} \textbf{Tauri 2.0} & Aplicaciones de escritorio ligeras y seguras \\[0.8em]
    \textcolor{isoGreen}{\faIcon{network-wired}} \textbf{gRPC} & Comunicaci\'on entre servicios de alto rendimiento \\[0.8em]
    \textcolor{isoBlue}{\faIcon{bolt}} \textbf{WebSockets} & Tiempo real bidireccional para notificaciones \\[0.8em]
\end{tabular}

\begin{cajadestacada}{Nota de Rendimiento}
    La combinaci\'on de Go para el backend y Tauri para el frontend proporciona un equilibrio \'optimo entre rendimiento y consumo de recursos, resultando en aplicaciones que pueden ejecutarse en hardware modesto.
\end{cajadestacada}

\clearpage
